% 原书 p.330–p.334（PDF p.338–p.342）

% >>> p338
\section{Arnold--Givental 猜想（Arnold--Givental's Conjecture）}

\begin{keytakeaways}
\begin{itemize}
  \item 几何图像：设辛流形 $(X,\omega)$ 带有一个反辛对合 $\tau$（即
        $\tau^*\omega=-\omega$），其不动轨迹 $\mathrm{Fix}(\tau)$ 是拉格朗日
        子流形。Arnold--Givental 猜想断言：任何紧支撑哈密顿同痕的时间 $t$ 映射
        都把 $\mathrm{Fix}(\tau)$ 移得与自身相交。
  \item 数学工具：$T^*M$ 上的齐次提升（把 $T^*M$ 的哈密顿同痕提升为
        $\dot T^*(M\times\R)$ 上的齐次哈密顿同痕）、GKS 层量子化
        $\mathcal K_\Phi$、层复合与 $SS$ 的函子性 (10.51)。
  \item 关键定理：定理 10.8.1（余切丛情形）与定理 10.8.2（微局部莫尔斯定理）；
        引理 10.8.1（$SS$ 只含非正余向量时层上同调「传播」）、命题 10.8.1
        （形变不变性）。
  \item 证明策略：用 $SS(\mathcal K_\Phi\circ\mathbf{k}_N)\subset
        \varphi_H^{-t}(SS(\mathbf{k}_N)\cap\dot T^*X)\cup0_X$ 把「层复合」
        化为「拉格朗日对应」，再用微局部莫尔斯定理与 $Rf_*$ 的前推公式把
        $H^*(M\times\R;\mathbf{k}_N)\ne0$ 转译为 $\varphi_H^{-t}(0_M)\cap0_M
        \ne\emptyset$。
\end{itemize}
\end{keytakeaways}

设 $(X^{2n},\omega)$ 是配备反辛对合 $\tau\colon X\to X$ 的辛流形，此处的
「反辛」意指 $\tau^*\omega=-\omega$。假设不动轨迹
$\mathrm{Fix}(\tau)\subset X$ 是拉格朗日子流形。下述例子展示了这一设定的两种
情形。

\begin{wexam}{10.8.1}
  (i) 对 $(X,\omega)=(T^*M,\omega_{\mathrm{can}})$，考虑由
  $\tau(x,p)=(x,-p)$（即把余向量变号）定义的映射
  $\tau\colon(X,\omega)\to(X,\omega)$。则
  \[
    \mathrm{Fix}(\tau)=0_M,
  \]
  它是 $(T^*M,\omega_{\mathrm{can}})$ 中的拉格朗日子流形。
  (ii) 对 $(X,\omega)=(\CP^n,\omega_{\mathrm{FS}})$，考虑由
  $\tau(z)=\bar z$ 定义的
  $\tau\colon(\CP^n,\omega_{\mathrm{FS}})\to(\CP^n,\omega_{\mathrm{FS}})$。则
  \[
    \mathrm{Fix}(\tau)=\RP^n,
  \]
  它也是 $(\CP^n,\omega_{\mathrm{FS}})$ 中的拉格朗日子流形。\hfill$\square$
\end{wexam}

下述猜想原来就是 [94] 中的猜想 C，它用（拉格朗日）子流形来描述哈密顿动力学
（更确切地说，哈密顿同痕）的刚性。相关讨论见 5.3 节。

% >>> p339
\begin{wconj}{10.8.1}
  设 $(X,\omega)$ 是配备反辛对合 $\tau\colon(X,\omega)\to(X,\omega)$ 的辛流形。
  假设不动轨迹 $\mathrm{Fix}(\tau)$ 是拉格朗日子流形，则对任何紧支撑哈密顿
  同痕 $\{\varphi_H^t\}_{t\in[0,1]}\colon[0,1]\times X\to X$，对任何
  $t\in[0,1]$ 都有
  \[
    (\varphi_H^t)^{-1}\bigl(\mathrm{Fix}(\tau)\bigr)\cap\mathrm{Fix}(\tau)
    \ne\emptyset.
  \]
\end{wconj}

该猜想已在多种情形下得到确认，包括例 10.8.1 中的两种情形。本节中，我们将基于
层量子化方法，对例 10.8.1 中 (i) 的情形 $(T^*M,\omega_{\mathrm{can}})$（其中
$\mathrm{Fix}(\tau)=0_M$）给出猜想 5.3.2 的一个证明。与推论 10.1.1（它是在
假设定理 10.1.1 的前提下与下述结果相同的结论）不同，定理 10.8.1 的证明完全
基于前几节所建立的层论机器。

% [Tutor Note: 精确陈述——猜想 10.8.1 的三个条件缺一不可：(i) $\tau$ 必须是
% \emph{反辛}对合 $\tau^*\omega=-\omega$（而非辛对合），这保证
% $\mathrm{Fix}(\tau)$ 的切空间是拉格朗日的（$T_xX=T_x\mathrm{Fix}\oplus
% \mathrm d\tau(T_x\mathrm{Fix})$ 且 $\omega$ 在 $\mathrm d\tau$ 下变号）；
% (ii) $\{\varphi_H^t\}$ 必须\emph{紧支撑}，否则平移反例（见注 10.8.1 (i)）
% 使结论不成立；(iii) 结论是「$\varphi_H^t$ 把不动轨迹移得必然与自身相交」，
% 它比 Arnold 猜想（哈密顿微分同胚的不动点）更强，因为 $\mathrm{Fix}(\tau)$
% 只是 $X$ 中的一个拉格朗日子流形而非整个流形。]

\begin{wthm}{10.8.1（余切丛上的 Arnold--Givental 猜想）}
  给定紧支撑哈密顿同痕
  $\Phi=\{\varphi_H^t\}_{t\in I}\colon I\times T^*M\to T^*M$，其中
  $I\subset\R$ 是含 $0\in\R$ 的区间，则对任何 $t\in I$ 都有
  \[
    (\varphi_H^t)^{-1}(0_M)\cap0_M\ne\emptyset.
  \]
\end{wthm}

\begin{wrem}{10.8.1}
  (i) 对一般的辛同胚，定理 10.8.1 的结论当然是错的。例如取 $M=S^1$，令
  $\varphi\colon T^*S^1\to T^*S^1$ 为纤维/竖直方向上的平移
  $\varphi(q,p)=(q,p+c)$（再乘上一个截断函数使 $\varphi$ 具有紧支撑）。则
  $\varphi$ 是辛同胚，但 $\varphi(0_{S^1})\cap0_{S^1}=\emptyset$。
  (ii) 在定理 10.8.1 中，若进一步假设
  $(\varphi_H^t)^{-1}(0_M)\pitchfork0_M$，则我们有更精细的估计（参看
  猜想 5.3.1 与猜想 5.3.2）
  \[
    \#\Bigl((\varphi_H^t)^{-1}(0_M)\cap0_M\Bigr)\ge\sum_i b_i(M;\mathbf{k}),
  \]
  其中 $b_i(M;\mathbf{k})$ 表示 $M$ 以 $\mathbf{k}$ 为系数的第 $i$ 个 Betti 数。
\end{wrem}

定理 10.8.1 证明中的关键步骤是下述引理，它看起来相当初等，但其证明却绝非
平凡（见 [135] 中的推论 5.4.19）。

\begin{wlem}{10.8.1}
  设 $\mathcal F\in\Op{Ob}(D(\Op{Sh}(\mathbf{k}_{\R})))$ 使得
  $SS(\mathcal F)|_{[a,b]}$ 只含非正余向量（对某闭区间 $[a,b]\subset\R$）。则
  \begin{equation}\tag{10.52}
    H^*\bigl((-\infty,a);\mathcal F\bigr)
    \simeq H^*\bigl((-\infty,b);\mathcal F\bigr).
  \end{equation}
\end{wlem}

换言之，层上同调群从 $a$「传播」到 $b$。

% [Tutor Note: 含糊点拨——引理 10.8.1 是 [135] 中推论 5.4.19 的特例，作者未给
% 证明。其直观是：$SS(\mathcal F)|_{[a,b]}$ 只含 $\eta\le0$ 的余向量，说明
% 在区间 $[a,b]$ 上方没有「正方向」的奇性传播，故 $\Gamma((-\infty,t);\mathcal
% F)$ 作为 $t$ 的函数在 $[a,b]$ 上是局部常值的，取导出函子后即得 (10.52)。
% 注意条件必须限制在 $[a,b]$ 上（即 $SS(\mathcal F)\cap T^*\R|_{[a,b]}$），
% 而不是整个 $SS(\mathcal F)$；否则结论一般不成立。此处「propagates from
% $a$ to $b$」是 10.6 节传播概念在一维参数族上的应用。]

下面是支持该引理的一个例子。

% >>> p340
\begin{wexam}{10.8.2}
  设 $\mathcal F=\mathbf{k}_{(0,1]}\in\Op{Ob}(\Op{Sh}(\mathbf{k}_\R))$。由
  例 10.6.1 与例 10.6.2 中的先前计算，我们知道 $SS(\mathcal F)$ 只含非正
  余向量（见图 10.7）。同时，我们断言
  \[
    H^0\bigl((-\infty,0);\mathcal F\bigr)
    \simeq H^0\bigl((-\infty,1);\mathcal F\bigr)(=\mathbf{k}).
  \]
  的确，对 (10.30) 作用 $\Gamma((-\infty,0))$，我们有
  \[
    0\to H^0\bigl((-\infty,0);\mathbf{k}_{(0,1]}\bigr)
    \to H^0\bigl((-\infty,0);\mathbf{k}_{(-\infty,1]}\bigr)
    \to H^0\bigl((-\infty,0);\mathbf{k}_{(-\infty,0]}\bigr)\to\cdots
  \]
  即为
  $0\to H^0((-\infty,0);\mathbf{k}_{(0,1]})\to\mathbf{k}
  \xrightarrow{\ \mathbf{1}_{\mathbf{k}}\ }\mathbf{k}\to\cdots$。因此
  $H^0((-\infty,0);\mathbf{k}_{(0,1]})=0$。类似地，可以验证
  $H^0((-\infty,1);\mathbf{k}_{(0,1]})=0$，因此引理 10.8.1 中的关系 (10.52)
  得到确认。\hfill$\square$
\end{wexam}

% [Tutor Note: 原文疑误——本例先说
% $H^0((-\infty,0);\mathcal F)\simeq H^0((-\infty,1);\mathcal F)(=\mathbf{k})$，
% 但紧接着的计算却得出两端均为 $0$（不是 $\mathbf{k}$）。按上下文，
% 「$(=\mathbf{k})$」应是笔误：对 $\mathcal F=\mathbf{k}_{(0,1]}$ 而言两端确实
% 都等于 $0$，(10.52) 的同构仍成立。此处保留原文写法并已记入审校记录。]

引理 10.8.1 推出下述相交结果，它本身也颇有意思。

\begin{wthm}{10.8.2（微局部莫尔斯定理）}
  设 $\mathcal F\in\Op{Ob}(D(\Op{Sh}(\mathbf{k}_X)))$，其中 $X$ 是闭的（或更
  一般地，$\Op{supp}(\mathcal F)$ 是紧的），$f\colon X\to\R$ 是可微函数。
  假设 $H^*(X;\mathcal F)\ne0$，则
  \[
    SS(\mathcal F)\cap\operatorname{graph}(\mathrm df)\ne\emptyset.
  \]
\end{wthm}

% [Tutor Note: 微局部刻画——定理 10.8.2 把经典莫尔斯理论中「非零同调迫使
% 临界点存在」的论断翻译为奇性支撑语言：$\operatorname{graph}(\mathrm df)
% \subset T^*X$ 与 $SS(\mathcal F)$ 相交，等价于说 $\mathcal F$ 沿
% $\mathrm df$ 方向必在某点\emph{不}传播；而这一点正是 $f$ 的临界点。注意
% 「$X$ 闭或 $\Op{supp}(\mathcal F)$ 紧」保证 $H^*(X;\mathcal F)$ 与紧支撑
% 前推的论证可以互译。该定理由 [135] 中的微局部莫尔斯理论给出，本书用引理
% 10.8.1 给出初等证明。]

\begin{wexam}{10.8.3}
  设 $\mathcal F=\mathbf{k}_X$，则 $H^*(X;\mathcal F)=H^*(X;\mathbf{k})\ne0$
  自动成立。同时，结论平凡地成立，因为 $f$ 总有临界点使 $\mathrm df(x)=0$，
  因此
  \[
    SS(\mathcal F)\cap\operatorname{graph}(\mathrm df)
    =SS(\mathbf{k}_X)\cap\operatorname{graph}(\mathrm df)
    =0_X\cap\operatorname{graph}(\mathrm df)\ne\emptyset,
  \]
  满足定理 10.8.2。\hfill$\square$
\end{wexam}

\begin{wproof}[定理 10.8.2 的证明（假设引理 10.8.1）]
  我们证明其逆否命题：若 $SS(\mathcal F)\cap\operatorname{graph}(\mathrm df)
  =\emptyset$，则 $H^*(X;\mathcal F)=0$。由前推的定义，我们有
  \[
    (f_*\mathcal F)\bigl((-\infty,a)\bigr)
    =\mathcal F\bigl(f^{-1}((-\infty,a))\bigr).
  \]
  换言之，
  $\Gamma((-\infty,a);f_*\mathcal F)
  =\Gamma(f^{-1}((-\infty,a));\mathcal F)$，这意味着
  \begin{equation}\tag{10.53}
    H^*\bigl((-\infty,a);Rf_*\mathcal F\bigr)
    =H^*\bigl(f^{-1}((-\infty,a));\mathcal F\bigr).
  \end{equation}
  回忆命题 10.7.2 中奇性支撑的前推公式，
  \[
    SS(Rf_*\mathcal F)\subset\bigl\{(y,\eta)\in T^*\R\ \big|\
    \exists(x,p)\in SS(\mathcal F)\ \text{使得}\ f(x)=y,\ f^*(\eta)=p\bigr\}.
  \]
  按定义 $f^*(\eta)=\eta\mathrm df$，其中 $\eta$ 被认同为一个数。则我们的
  假设（与 $\operatorname{graph}(\mathrm df)$ 不相交）特别地意味着
  $\eta\ne1$。此外，对任何 $\eta>0$，
  $(x,\eta\mathrm df)\notin SS(\mathcal F)$，因为 $SS(\mathcal F)$ 是锥的。
  由此我们断定
  \[
    SS(Rf_*\mathcal F)\cap T^*\R|_{[a,b]}
    \subset\bigl\{(y,\eta)\in T^*\R\mid\eta\le0\bigr\}
  \]

% >>> p341
  对任何区间 $[a,b]\subset\R$。因此由引理 10.8.1，
  \[
    H^*\bigl((-\infty,a);Rf_*\mathcal F\bigr)
    \simeq H^*\bigl((-\infty,b);Rf_*\mathcal F\bigr).
  \]
  结合 (10.53)，上述同构转化为
  $H^*(f^{-1}((-\infty,a));\mathcal F)
  \simeq H^*(f^{-1}((-\infty,b));\mathcal F)$。正如莫尔斯理论中那样，我们将
  考虑「过滤的」上同调群。这里有两个极端情形。
  \begin{enumerate}[label=(\roman*)]
    \item 当 $a$ 充分负时，下水平集 $f^{-1}((-\infty,a))=\emptyset$，这意味
          着 $H^*(f^{-1}((-\infty,a));\mathcal F)=0$。
    \item 当 $b$ 充分正时，下水平集 $f^{-1}((-\infty,b))=X$，这意味着
          $H^*(f^{-1}((-\infty,b));\mathcal F)=H^*(X;\mathcal F)$。
  \end{enumerate}
  结合 (i) 与 (ii)，得到 $H^*(X;\mathcal F)=0$，这就完成了证明。
\end{wproof}

我们差不多可以给出定理 10.8.1 的证明了。在此之前，还需要两个准备结果。第一个
结果说明何时可以在不改变所得层上同调群的前提下「形变」层 $\mathcal F$。第二个
结果是定理 10.8.2 的一个显式应用。

\begin{wprop}{10.8.1}
  对 $\mathcal F\in\Op{Ob}(D(\Op{Sh}(\mathbf{k}_{X\times I})))$，若
  $SS(\mathcal F)\cap(0_X\times T^*I)\subset0_{X\times I}$，则可以如下形变
  层的系数：
  \[
    H^*(X;\mathcal F|_{t=0})\simeq H^*(X;\mathcal F|_{t=1}).
  \]
\end{wprop}

\begin{wproof}
  记 $q\colon X\times I\to I$。回忆命题 10.7.2 中奇性支撑的前推公式，
  \[
    SS(Rq_*\mathcal F)\subset\bigl\{(t,\eta)\in T^*I\ \big|\
    \exists(x,t)\in X\times I\ \text{使得}\ (x,t,0,\eta)\in SS(\mathcal F)\bigr\}.
  \]
  关于 $SS(\mathcal F)$ 的假设说明 $\eta=0$，故 $SS(Rq_*\mathcal F)\subset0_I$。
  这意味着 $Rq_*\mathcal F$ 是常值层（依据 [135] 中的习题 II.4）。这又意味着
  $R\Gamma(X,\mathcal F|_{t=0})=R\Gamma(X,\mathcal F|_{t=1})$。由于
  $\mathcal F$ 有紧支撑，这等价于 $(Rq_*\mathcal F)_0=(Rq_*\mathcal F)_1$。
\end{wproof}

% [Tutor Note: 原文疑误——命题 10.8.1 的证明中，作者写
% "$SS(Rq_*\mathcal F)\subset0_{X\times I}$"，但 $Rq_*\mathcal F$ 是 $I$ 上的
% 层，其奇性支撑应落在 $T^*I$ 中，故正确写法应为
% $SS(Rq_*\mathcal F)\subset0_I$。此处保留原文写法并已记入审校记录。]

下面是命题 10.8.1 的一个具体应用场景。对齐次哈密顿同痕
$\Phi=\{\varphi_H^t\}_{t\in I}\colon I\times\dot T^*X\to\dot T^*X$，由定理
10.6.1 与例 10.7.5，存在
$\mathcal K_\Phi\in\Op{Ob}(D(\Op{Sh}(\mathbf{k}_{X\times X\times I})))$，
使得对任何 $\mathcal F\in\Op{Ob}(D(\Op{Sh}(\mathbf{k}_X)))$，我们有
$\mathcal K_\Phi\circ\mathcal F
\in\Op{Ob}(D(\Op{Sh}(\mathbf{k}_{X\times I})))$，并且
\[
  SS(\mathcal K_\Phi\circ\mathcal F)\subset
  \Bigl(\bigcup_{t\in I}\varphi_H^{-t}\bigl(SS(\mathcal F)\cap\dot T^*X\bigr)
  \cup0_X\Bigr)\times0_I.
\]
这意味着
$SS(\mathcal K_\Phi\circ\mathcal F)\cap(0_X\times T^*I)\subset0_{X\times I}$。
因此由上面的命题 10.8.1，
\[
  H^*\bigl(X;(\mathcal K_\Phi\circ\mathcal F)|_{t=0}\bigr)
  \simeq H^*\bigl(X;(\mathcal K_\Phi\circ\mathcal F)|_{t=1}\bigr).
\]

% >>> p342
最后，由定理 10.6.1 中的 (ii) 我们有
\[
  \begin{aligned}
    H^*\bigl(X;(\mathcal K_\Phi\circ\mathcal F)|_{t=0}\bigr)
    &=H^*\bigl(X;(\mathcal K_\Phi|_{t=0}\circ\mathcal F)\bigr)\\
    &=H^*\bigl(X;\mathbf{k}_\Delta\circ\mathcal F\bigr)
    =H^*(X;\mathcal F).
  \end{aligned}
\]
我们断定
$H^*(X;\mathcal F)\simeq H^*(X;\mathcal K_\Phi|_{t=1}\circ\mathcal F)$。
事实上，通过重新参数化时间 $t$，我们有
\begin{equation}\tag{10.54}
  H^*(X;\mathcal F)\simeq H^*(X;\mathcal K_\Phi|_t\circ\mathcal F)
\end{equation}
对任何 $t\in I$。结合定理 10.8.2，这直接推出下述结果。详细留给读者作为练习
（提示：利用 (10.51) 与 (10.54)）。

\begin{wprop}{10.8.2（练习）}
  设 $\Phi=\{\varphi_H^t\}_{t\in I}\colon I\times\dot T^*X\to\dot T^*X$ 是
  齐次哈密顿同痕，$f\colon X\to\R$ 是满足对任何 $x\in X$ 有
  $\mathrm df(x)\ne0$ 的函数。则对任何 $t\in I$，若
  $\mathcal F\in\Op{Ob}(D(\Op{Sh}(\mathbf{k}_X)))$ 有紧支撑，则
  \begin{equation}\tag{10.55}
    (\varphi_H^{-t})\bigl(SS(\mathcal F)\cap\dot T^*X\bigr)
    \cap\operatorname{graph}(\mathrm df)\ne\emptyset.
  \end{equation}
\end{wprop}

最后，这里是定理 10.8.1 的证明。

\begin{wproof}[定理 10.8.1 的证明]
  设 $\Phi=\{\varphi_H^t\}_{t\in I}\colon I\times T^*M\to T^*M$ 是哈密顿
  同痕。回忆定理 10.6.1 之前的讨论，我们可以把它提升为齐次哈密顿同痕
  $\widetilde\Phi=\{\varphi_{\widetilde H}^t\}\colon
  I\times\dot T^*(M\times\R)\to\dot T^*(M\times\R)$，具体定义为
  \[
    \bigl(t,(m,p),(s,\sigma)\bigr)
    \longmapsto
    \Bigl(\sigma\cdot\varphi_H^t(m,p/\sigma),\ 
    \bigl(s+u(t,(m,p/\sigma)),\sigma\bigr)\Bigr).
  \]
  选取 $N=M\times\{0\}\subset M\times\R$ 并考虑
  $\mathcal F=\mathbf{k}_N$。应用定理 10.5.1 中的复合公式（在例 10.5.5 末尾），
  我们有
  \[
    H^*(M\times\R;\mathbf{k}_N)=H^*(N;\mathbf{k})\ne0.
  \]
  考虑由 $f(m,s)=s$ 定义的投影 $f\colon M\times\R\to\R$。由于 $\mathrm df$
  在 $M\times\R$ 上恒不为零，命题 10.8.2 适用。特别地，我们可以显式写出
  (10.55) 中的各分量如下：
  \[
    \begin{aligned}
      (\varphi_{\widetilde H}^{-t})\bigl(SS(\mathcal F)\cap\dot T^*X\bigr)
      &=(\varphi_{\widetilde H}^{-t})\bigl(\nu_N^*(M\times\R)
        \cap\dot T^*(M\times\R)\bigr)\\
      &=(\varphi_{\widetilde H}^{-t})\bigl(\{(m,0,s,\sigma)\mid
        m\in M,\ s\in\R,\ \sigma\ne0\}\bigr)\\
      &=\bigl\{\bigl(\sigma\cdot\varphi_H^{-t}(m,0),
        (s+u(t,(m,0)),\sigma)\bigr)\ \big|\
        m\in M,\ s\in\R,\ \sigma\ne0\bigr\}.
    \end{aligned}
  \]
  同时，
  $\operatorname{graph}(\mathrm df)
  =\bigl\{((m',0),(s',1))\in T^*M\times T^*\R\mid m'\in M,\ s'\in\R\bigr\}$。
  由于 $\operatorname{graph}(\mathrm df)$ 与
  $(\varphi_{\widetilde H}^{-t})(SS(\mathcal F)\cap\dot T^*X)$ 的交非空，取
  $\sigma=1$ 便知存在 $m,m'\in M$ 使得
  \[
    (\varphi_H^{-t})(m,0)=(m',0).
  \]
  换言之，$(\varphi_H^{-t})(0_M)\cap0_M$ 非空。这就完成了证明。
\end{wproof}

% [Tutor Note: 层论证明策略（全节要点小结）——定理 10.8.1 的证明链条可概括为
% 五步：(1) 把 $T^*M$ 上的哈密顿同痕齐次提升到 $\dot T^*(M\times\R)$，
% 从而进入 GKS 层量子化的适用范围（定理 10.6.1）；(2) 取 $N=M\times\{0\}$、
% $\mathcal F=\mathbf{k}_N$，其层上同调 $H^*(N;\mathbf{k})\ne0$；(3) 用层复合
% 与 (10.51) 把 $SS(\mathcal K_\Phi\circ\mathcal F)$ 控制在
% $\varphi_H^{-t}(SS(\mathcal F)\cap\dot T^*X)\cup0_X$ 的竖直平移内，
% 并用命题 10.8.1（形变不变性）说明 $t=0$ 与 $t=1$ 的上同调相同；
% (4) 把 $\mathcal F=\mathbf{k}_N$ 的 $SS$ 显式算出（$\nu_N^*(M\times\R)$），
% 用微局部莫尔斯定理（定理 10.8.2）作用于 $f(m,s)=s$，得到
% (10.55) 的非空相交；(5) 在交集中取 $\sigma=1$，读出
% $(\varphi_H^{-t})(m,0)=(m',0)$。全程未使用任何伪全纯曲线，这正是层量子化
% 方法相对于 Floer 理论的优势所在。]

% [Tutor Note: 定理编号与交叉引用提示——本节引用的 10.x 编号均照原书：
% 定理 10.6.1（GKS 层量子化）、定理 10.5.1（复合公式）、例 10.5.5、
% 例 10.6.1--10.6.5、命题 10.7.2/10.7.3、引理 10.8.1、命题 10.8.1、
% 定理 10.8.1/10.8.2、推论 10.1.1，以及 5.3 节的猜想 5.3.1、5.3.2。
% 校样时请逐一核对上述编号与原文的对应关系（尤其不要与 10.8.1 例混淆）。]
